\chapter{Метод переноса динамических представлений и методика сравнительной оценки}

\section{Формирование временного представления нейродинамики}
\ResultPending{Будут определены сопоставимые характеристики динамики, размерность и временной масштаб представления.}

\section{Алгоритм предобучения и адаптации представления}
\ResultPending{Планируется переносить начальные параметры временного модуля с адаптацией по обучающим данным человека.}

\section{Сопряжение динамического представления с физическим оператором ECAP}
\ResultPending{Будет определено сопряжение временного модуля с моделью рекрутирования волокон и отдельным оператором регистрации сигнала.}
Планируемое сопряжение компонентов показано на иллюстрации «Рисунок~\ref{fig:planned-method}». Предполагается перенос начальных параметров временного модуля; анатомическое соответствие видов не предполагается.
\begin{figure}[htbp]
\centering
\begin{tikzpicture}[node distance=8mm]
\node[research process] (graph) {Коннектом\\Drosophila};
\node[research process,right=of graph] (dynamics) {Нейродинамика\\и временные ряды};
\node[research process,right=of dynamics] (pretraining) {Предобучение\\временного модуля};
\draw[research flow] (graph)--(dynamics);
\draw[research flow] (dynamics)--(pretraining);
\node[research process,below=29mm of pretraining] (adaptation) {Адаптация\\временного модуля};
\node[research process,left=of adaptation] (recruitment) {Рекрутирование\\волокон человека};
\node[research process,left=of recruitment] (recording) {Регистрация сигнала\\и кривая роста};
\draw[research flow] (pretraining)--node[right,text width=24mm,align=left,font=\fontsize{13}{15.7}\selectfont]{Начальные\\параметры}(adaptation);
\draw[research flow] (adaptation)--(recruitment);
\draw[research flow] (recruitment)--(recording);
\node[draw,dashed,fit=(graph)(dynamics)(pretraining),inner sep=2mm,line width=0.5pt] (source) {};
\node[anchor=south west,font=\fontsize{13}{15.7}\selectfont] at (source.north west) {Предобучение на Drosophila};
\node[draw,dashed,fit=(adaptation)(recruitment)(recording),inner sep=2mm,line width=0.5pt] (target) {};
\node[anchor=south west,font=\fontsize{13}{15.7}\selectfont] at (target.north west) {Целевая модель человеческого сигнала};
\end{tikzpicture}

\caption{Схема планируемого метода моделирования ECAP}
\label{fig:planned-method}
\end{figure}
\FloatBarrier

\section{Постановка задачи моделирования кривой роста ECAP и выбор критериев оценки}
\ResultPending{Ошибка амплитуды кривой роста будет рассчитана для каждого участника и усреднена между участниками. Сравнение будет включать обучение без переноса, предобучение на биофизической модели волокон человека, контроль временной структуры и топологии, а также простую модель кривой роста.}
Планируемые сравнения приведены в материале «Таблица~\ref{tab:planned-controls}»; их предполагается проводить при сопоставимых данных, целевой архитектуре, физическом операторе и вычислительных затратах.
\begin{table}[htbp]
\caption{Планируемые контрольные сравнения}
\label{tab:planned-controls}
\centering
\ResearchTableFont
\renewcommand{\arraystretch}{1.15}
\begin{tabularx}{\linewidth}{|L{0.43\linewidth}|Y|}
\hline Контроль & Проверяемый вклад \\\hline
Обучение только на человеческих ECAP & Дополнительная ценность предобучения \\\hline
Предобучение на модели волокон человека & Выбор источника предобучения \\\hline
Перемешанная временная динамика & Временная структура активности \\\hline
Перестроенные коннектомы & Топология при сохранении заданных характеристик графа \\\hline
Простая модель кривой роста & Необходимость временного модуля \\\hline
\end{tabularx}
\end{table}
\FloatBarrier

\section{Программная реализация метода}
\ResultPending{Будут описаны программная реализация, происхождение данных и параметры, необходимые для воспроизведения расчётов.}

\section{Формирование независимых выборок и план статистического анализа}
\ResultPending{Допуски, объём выборки, правила разбиения и калибровки будут зафиксированы до доступа к независимой проверочной выборке.}

\section{Анализ неопределённости и влияния параметров измерения}
\ResultPending{Планируется исследовать чувствительность к геометрии, параметрам стимуляции и условиям регистрации, а также калибровку прогнозных интервалов.}

\section{Область применимости метода и условия получения достоверных оценок}
\ResultPending{Будут установлены условия применимости модели и границы допустимых технических выводов.}

\section*{Выводы по главе}
\ResultPending{Выводы будут сформулированы после реализации задач главы.}
